% Chapter 2 draft for textbook inclusion. Assumes a textbook preamble with amsmath and amssymb.
\chapter{Prediction Error, Flexibility, and Generalization}
\label{chap:prediction-error-flexibility-generalization}

\CoreCompress{
\section{Why Prediction Error Matters}

The previous chapter introduced the supervised learning model
\begin{equation}
  Y = f_0(\mathbf{X}) + \epsilon,
  \label{eq:ch2-true-model}
\end{equation}
where $\mathbf{X}=(X_1,\ldots,X_p)^T$ is the vector of input variables, $Y$ is the response, $f_0$ is the true systematic relationship, and $\epsilon$ is the error term.  In this chapter we ask a prediction-focused question:
\begin{quote}
  Given the information in $\mathbf{X}$, how accurately can we predict $Y$?
\end{quote}
This question makes precise the phrase ``all models are wrong.''  A model is not judged by whether it perfectly describes reality.  It is judged by how much useful signal it captures, how much error remains, and whether its performance holds up on data it has not already seen.

To avoid a common notation problem, this chapter uses $f_0$ for the true unknown function and $\hat f$ for a fitted model.  The true function $f_0$ is fixed but unknown.  The fitted function $\hat f$ is random because it depends on the training data used to fit it.  If a different training data set were collected, the fitted model would generally be different.

\section{Mean Squared Error}

The most common measure of prediction error for a quantitative response is \emph{mean squared error}.  For a new observation with predictors $\mathbf{x}_0$ and response $Y_0$, the prediction is
\[
  \widehat{Y}_0 = \hat f(\mathbf{x}_0).
\]
The squared prediction error is
\[
  \left(Y_0 - \hat f(\mathbf{x}_0)\right)^2.
\]
The mean squared error, or MSE, is the expected squared prediction error:
\begin{equation}
  \operatorname{MSE}(\mathbf{x}_0)
  = \mathbb{E}\left[\left(Y_0 - \hat f(\mathbf{x}_0)\right)^2\right].
  \label{eq:ch2-mse-fixed-x}
\end{equation}
The expectation in Equation~\eqref{eq:ch2-mse-fixed-x} averages over the randomness in the training data used to produce $\hat f$ and the randomness in the new response $Y_0$.

For an observed data set, we often compute an empirical version of MSE.  Given observations $(\mathbf{x}_1,y_1),\ldots,(\mathbf{x}_n,y_n)$ and predictions $\hat y_i=\hat f(\mathbf{x}_i)$, the empirical MSE is
\begin{equation}
  \operatorname{MSE}
  = \frac{1}{n}\sum_{i=1}^n (y_i - \hat y_i)^2.
  \label{eq:ch2-empirical-mse}
\end{equation}
Large errors are penalized heavily because the errors are squared.  This makes MSE mathematically convenient, but it also means that a few large mistakes can dominate the score.

The square root of MSE is called \emph{root mean squared error}, or RMSE:
\begin{equation}
  \operatorname{RMSE}
  = \sqrt{\operatorname{MSE}}.
  \label{eq:ch2-rmse}
\end{equation}
RMSE is often easier to interpret because it is measured in the same units as the response.  For example, if $Y$ is measured in hours, RMSE is also measured in hours; if $Y$ is measured in thousands of dollars, RMSE is measured in thousands of dollars.

\paragraph{Math Foundations Connection.}
Let $\mathbf{y}=(y_1,\ldots,y_n)^T$ be the vector of observed responses and let $\hat{\mathbf{y}}=(\hat y_1,\ldots,\hat y_n)^T$ be the vector of predictions.  The residual vector is
\[
  \mathbf{e}=\mathbf{y}-\hat{\mathbf{y}}.
\]
Using the Euclidean norm,
\[
  \|\mathbf{e}\| = \sqrt{e_1^2+\cdots+e_n^2},
\]
the empirical MSE can be written as
\begin{equation}
  \operatorname{MSE}=\frac{1}{n}\|\mathbf{y}-\hat{\mathbf{y}}\|^2.
  \label{eq:ch2-mse-norm}
\end{equation}
Thus MSE is a scaled squared distance between the observed response vector and the fitted response vector.  This geometric view will become especially useful when we later interpret least squares regression as a projection problem.  See the Math Foundations textbook, Chapters 2, 3, 5, and 12.

\section{The Bias--Variance Decomposition}

The bias--variance decomposition explains why prediction is limited even when we use a good model.  It separates expected prediction error into three parts:
\begin{equation}
  \text{MSE} = \text{bias}^2 + \text{model variance} + \text{noise}.
  \label{eq:ch2-bv-summary}
\end{equation}
This equation is one of the central ideas in statistical learning.

\subsection{Set-Up for the Decomposition}

Fix a predictor value $\mathbf{x}_0$.  Suppose the new response is generated by
\begin{equation}
  Y_0=f_0(\mathbf{x}_0)+\epsilon_0,
  \label{eq:ch2-new-response}
\end{equation}
where
\begin{align*}
  \mathbb{E}[\epsilon_0] &= 0,\\
  \operatorname{Var}(\epsilon_0) &= \sigma^2,
\end{align*}
with $\epsilon_0$ independent of the training data used to construct $\hat f$.  We want the expected squared prediction error
\[
  \mathbb{E}\left[\left(Y_0-\hat f(\mathbf{x}_0)\right)^2\right].
\]
Substitute Equation~\eqref{eq:ch2-new-response}:
\begin{align}
  \mathbb{E}\left[\left(Y_0-\hat f(\mathbf{x}_0)\right)^2\right]
  &= \mathbb{E}\left[\left(f_0(\mathbf{x}_0)+\epsilon_0-\hat f(\mathbf{x}_0)\right)^2\right].
  \label{eq:ch2-substitution}
\end{align}
Rearrange the expression inside the square:
\[
  f_0(\mathbf{x}_0)+\epsilon_0-\hat f(\mathbf{x}_0)
  = \left(f_0(\mathbf{x}_0)-\hat f(\mathbf{x}_0)\right)+\epsilon_0.
\]
Now expand the square:
\begin{align*}
  &\mathbb{E}\left[\left(\left(f_0(\mathbf{x}_0)-\hat f(\mathbf{x}_0)\right)+\epsilon_0\right)^2\right] \\
  &\quad = \mathbb{E}\left[\left(f_0(\mathbf{x}_0)-\hat f(\mathbf{x}_0)\right)^2\right]
  +2\mathbb{E}\left[\epsilon_0\left(f_0(\mathbf{x}_0)-\hat f(\mathbf{x}_0)\right)\right]
  +\mathbb{E}[\epsilon_0^2].
\end{align*}
The middle term is zero because $\epsilon_0$ has mean zero and is independent of the fitted model.  Also, since $\mathbb{E}[\epsilon_0]=0$, we have $\mathbb{E}[\epsilon_0^2]=\operatorname{Var}(\epsilon_0)=\sigma^2$.  Therefore,
\begin{equation}
  \mathbb{E}\left[\left(Y_0-\hat f(\mathbf{x}_0)\right)^2\right]
  = \mathbb{E}\left[\left(f_0(\mathbf{x}_0)-\hat f(\mathbf{x}_0)\right)^2\right]
  + \sigma^2.
  \label{eq:ch2-mse-signal-plus-noise}
\end{equation}
The final term, $\sigma^2$, is the noise variance.  It is the variation in $Y$ that remains even if $f_0$ were known perfectly.

\subsection{Separating Bias and Variance}

We now decompose
\[
  \mathbb{E}\left[\left(f_0(\mathbf{x}_0)-\hat f(\mathbf{x}_0)\right)^2\right].
\]
Let
\[
  m(\mathbf{x}_0)=\mathbb{E}\left[\hat f(\mathbf{x}_0)\right]
\]
be the average prediction we would get at $\mathbf{x}_0$ over many possible training data sets.  Add and subtract $m(\mathbf{x}_0)$ inside the square:
\begin{align*}
  \hat f(\mathbf{x}_0)-f_0(\mathbf{x}_0)
  &= \left(\hat f(\mathbf{x}_0)-m(\mathbf{x}_0)\right)
  + \left(m(\mathbf{x}_0)-f_0(\mathbf{x}_0)\right).
\end{align*}
Squaring and taking expectations gives
\begin{align*}
  \mathbb{E}\left[\left(\hat f(\mathbf{x}_0)-f_0(\mathbf{x}_0)\right)^2\right]
  &= \mathbb{E}\left[\left(\hat f(\mathbf{x}_0)-m(\mathbf{x}_0)\right)^2\right]
  + \left(m(\mathbf{x}_0)-f_0(\mathbf{x}_0)\right)^2 \\
  &\quad + 2\left(m(\mathbf{x}_0)-f_0(\mathbf{x}_0)\right)
  \mathbb{E}\left[\hat f(\mathbf{x}_0)-m(\mathbf{x}_0)\right].
\end{align*}
The final expectation is zero by the definition of $m(\mathbf{x}_0)$, so the cross term vanishes.  Hence,
\begin{equation}
  \mathbb{E}\left[\left(\hat f(\mathbf{x}_0)-f_0(\mathbf{x}_0)\right)^2\right]
  = \operatorname{Var}\left(\hat f(\mathbf{x}_0)\right)
  + \operatorname{Bias}\left(\hat f(\mathbf{x}_0)\right)^2,
  \label{eq:ch2-bias-variance-no-noise}
\end{equation}
where
\begin{align}
  \operatorname{Var}\left(\hat f(\mathbf{x}_0)\right)
  &= \mathbb{E}\left[\left(\hat f(\mathbf{x}_0)-\mathbb{E}[\hat f(\mathbf{x}_0)]\right)^2\right],
  \label{eq:ch2-model-variance}\\
  \operatorname{Bias}\left(\hat f(\mathbf{x}_0)\right)
  &= \mathbb{E}\left[\hat f(\mathbf{x}_0)\right]-f_0(\mathbf{x}_0).
  \label{eq:ch2-model-bias}
\end{align}
Combining Equations~\eqref{eq:ch2-mse-signal-plus-noise} and~\eqref{eq:ch2-bias-variance-no-noise}, we obtain
\begin{equation}
  \boxed{
  \operatorname{MSE}(\mathbf{x}_0)
  = \operatorname{Bias}\left(\hat f(\mathbf{x}_0)\right)^2
  + \operatorname{Var}\left(\hat f(\mathbf{x}_0)\right)
  + \operatorname{Var}(\epsilon_0).}
  \label{eq:ch2-bias-variance-final}
\end{equation}
This is the bias--variance decomposition.

\subsection{Interpreting the Three Terms}

The \emph{squared bias} term measures systematic error.  It is large when the average fitted model misses the true function.  A model that is too simple can have high bias.  For example, if the true relationship is curved but we insist on using a straight line, the fitted line may miss the main shape of the data no matter how much data we collect.

The \emph{model variance} term measures how much the fitted model changes from one training data set to another.  A highly flexible model may fit one training sample very closely but produce a noticeably different function if the sample changes.  This instability is model variance.

The \emph{noise} term, $\operatorname{Var}(\epsilon_0)$, is irreducible error.  It is the part of the response that cannot be predicted from $\mathbf{X}$ using any method, because it comes from unmeasured factors, measurement error, randomness in the system, or other variation not captured by the inputs.

This distinction gives us the language of \emph{reducible} and \emph{irreducible} error.  Bias and model variance are reducible in principle: we can try different model classes, collect better data, improve features, or tune the model.  The noise term is irreducible with respect to the inputs we have chosen.  The best possible prediction error approaches this irreducible error from above; it cannot go below it when the goal is to predict the actual future response $Y$.

\section{Training Error and Test Error}

Prediction models are trained on data.  Let
\[
  \mathcal{D}_{\text{train}}
  = \{(\mathbf{x}_1,y_1),\ldots,(\mathbf{x}_{n_{\text{train}}},y_{n_{\text{train}}})\}
\]
be the training data used to fit $\hat f$.  The \emph{training MSE} is
\begin{equation}
  \operatorname{MSE}_{\text{train}}
  = \frac{1}{n_{\text{train}}}\sum_{(\mathbf{x}_i,y_i)\in \mathcal{D}_{\text{train}}}
  \left(y_i-\hat f(\mathbf{x}_i)\right)^2.
  \label{eq:ch2-training-mse}
\end{equation}
Training MSE measures how well the model performs on data it already used to learn.

The prediction problem, however, is not mainly about the training data.  If we train a model using data from existing diabetes patients, we already know which of those patients had diabetes.  The real goal is to make accurate predictions for future patients whose outcomes are not yet known.  In the same way, a logistics forecast, readiness forecast, or maintenance forecast is useful only if it generalizes beyond the cases already used to build it.

Let
\[
  \mathcal{D}_{\text{test}}
  = \{(\mathbf{x}_1,y_1),\ldots,(\mathbf{x}_{n_{\text{test}}},y_{n_{\text{test}}})\}
\]
be a separate set of observations not used to fit the model.  The \emph{test MSE} is
\begin{equation}
  \operatorname{MSE}_{\text{test}}
  = \frac{1}{n_{\text{test}}}\sum_{(\mathbf{x}_i,y_i)\in \mathcal{D}_{\text{test}}}
  \left(y_i-\hat f(\mathbf{x}_i)\right)^2.
  \label{eq:ch2-test-mse}
\end{equation}
Test MSE measures generalization: the model's ability to predict unseen data.

Training MSE is usually optimistic.  Many fitting procedures explicitly choose model parameters to reduce training error.  Therefore, a low training MSE does not guarantee a low test MSE.  The model may simply have learned the idiosyncrasies of the training sample.

\section{Flexibility}

A model's \emph{flexibility} is its capacity to approximate a wide range of functional shapes.  A model is more flexible if it can represent a larger set of possible relationships $f(\mathbf{x})$.

For example, a straight-line model in one predictor can represent only lines:
\[
  \hat f(x)=\hat\beta_0+\hat\beta_1x.
\]
A polynomial model with many terms can represent many more shapes:
\[
  \hat f(x)=\hat\beta_0+\hat\beta_1x+\hat\beta_2x^2+\cdots+\hat\beta_dx^d.
\]
A smoothing spline, decision tree, random forest, boosted model, support vector machine, or neural network may be more flexible still, depending on how it is tuned.

Flexibility is not automatically good or bad.  Its value depends on the relationship between the true function, the amount of noise, and the amount of data available.  More flexible methods can reduce bias because they can capture more complex structure.  At the same time, more flexible methods tend to increase model variance because they are more sensitive to the particular training sample.

As a general rule:
\begin{equation}
  \text{more flexibility} \quad \Longrightarrow \quad
  \text{lower bias but higher variance}.
  \label{eq:ch2-flexibility-rule}
\end{equation}
This is only a general rule, not a law that applies identically in every model class.  It is nevertheless a useful guide for thinking about prediction error.

\subsection{Degrees of Freedom as a Flexibility Measure}

For some models, flexibility can be summarized by a quantity called \emph{degrees of freedom}.  In smoothing splines, for example, a low-degree-of-freedom fit is smoother and less flexible, while a high-degree-of-freedom fit is more wiggly and more flexible.  In polynomial regression, the number of estimated parameters provides a simple proxy for flexibility.  A model with two parameters is less flexible than a model with thirteen parameters.

The important pattern is this: as flexibility increases, training MSE usually decreases, because the model has more freedom to chase the training observations.  Test MSE often decreases at first, but then may increase when the model becomes too flexible.

\section{Overfitting and Overtraining}

\emph{Overfitting} occurs when a model follows noise or accidental patterns in the training data rather than the underlying relationship.  The model looks impressive on the data used to fit it, but performs poorly on new observations.  In the lecture slides, this is also described as \emph{overtraining}: an overabundance of flexibility causes test MSE to increase compared with simpler models.

Overfitting is undesirable because prediction is about future or unseen data.  A flexible model may fit the training data so well that it learns artifacts specific to that sample.  These artifacts do not appear in future data, so test performance suffers.

A useful mental picture is the following.  Suppose the true relationship is a smooth curve.  A very simple model may underfit: it misses important curvature and has high bias.  A moderately flexible model may capture the main trend and generalize well.  A very flexible model may twist through individual noisy observations; it has low training error but high variance and worse test error.

Therefore, the goal is not to make training MSE as small as possible.  The goal is to choose a model whose predictions generalize.

\section{Train/Test Splits}

A train/test split is a simple way to estimate generalization performance.

\begin{enumerate}
  \item Randomly divide the available data into a training set and a test set.
  \item Fit the model using only the training set.
  \item Predict the responses in the test set.
  \item Compute test MSE or test RMSE using the test-set predictions.
  \item Compare candidate models using test performance, not training performance alone.
\end{enumerate}

The key rule is that the test data must remain unseen during model fitting.  If the test data are used to choose transformations, tune flexibility, select predictors, or adjust the model, then they are no longer a clean test of generalization.  Later in the course we will use cross-validation to improve on a single train/test split, but the principle is the same: separate the data used to fit the model from the data used to evaluate the model.

\subsection{A Simulation Example}

The lecture simulation uses the model
\begin{equation}
  Y=f_0(X)+\epsilon,
  \label{eq:ch2-simulation-model}
\end{equation}
where
\begin{align*}
  f_0(X) &= X^3,\\
  X &\sim \operatorname{Uniform}(-3,3),\\
  \epsilon &\sim \mathcal{N}(0,9^2).
\end{align*}
The true function is cubic, but the observed data are noisy.  A train/test split separates the simulated observations into data used to fit the model and data used to evaluate it.

Consider four candidate models with increasing numbers of parameters.  The lecture example reports the following pattern:
\begin{center}
\begin{tabular}{c|c|c|c}
\hline
Model & Number of Parameters & Training MSE & Test MSE \\
\hline
1 & 2  & 101.99 & 111.82 \\
2 & 4  & 84.56  & 91.15  \\
3 & 9  & 83.52  & 90.60  \\
4 & 13 & 82.46  & 92.55  \\
\hline
\end{tabular}
\end{center}
The training MSE decreases as the number of parameters increases.  This is expected: more parameters give the fitted model more freedom to reduce training error.  The test MSE, however, decreases at first and then increases.  The 13-parameter model has the lowest training MSE, but not the lowest test MSE.  In this example, the 9-parameter model gives the best test performance among the four candidates.

The lesson is not that nine parameters are always best.  The lesson is that model selection is a balance between flexibility and test error.  A model that is too rigid underfits; a model that is too flexible overfits.  The best predictive model is usually somewhere between these extremes.

\section{Parametric and Nonparametric Estimators}

Another way to describe model classes is to distinguish \emph{parametric} and \emph{nonparametric} methods.

\subsection{Parametric Methods}

A method is parametric if the function $f$ is described using a finite number of parameters.  The model may be simple or complicated, linear or nonlinear, but it is governed by a finite parameter vector.

The standard linear model is parametric:
\begin{equation}
  f(\mathbf{X})
  = \beta_0 + \beta_1X_1+\cdots+\beta_pX_p
  = \beta_0+\boldsymbol{\beta}^T\mathbf{X},
  \label{eq:ch2-parametric-linear}
\end{equation}
where
\[
  \boldsymbol{\beta}=(\beta_1,\ldots,\beta_p)^T.
\]
Once we assume the form in Equation~\eqref{eq:ch2-parametric-linear}, learning $f$ becomes the problem of estimating the finite list of parameters $\beta_0,\beta_1,\ldots,\beta_p$.

\paragraph{Math Foundations Connection.}
The term $\boldsymbol{\beta}^T\mathbf{X}$ is an inner product.  It is also a linear combination of the entries of $\mathbf{X}$ with coefficients $\beta_1,\ldots,\beta_p$.  This is why linear models connect naturally to the vector notation from Math Foundations: the model turns a feature vector into a prediction by weighting its components and adding the weighted components together.  See the Math Foundations textbook, Chapters 2, 3, and 5.

Parametric does not mean ``good'' or ``simple.''  For example, one could write a finite-parameter model such as
\[
  f(\mathbf{X})
  = \beta_0 + \beta_{13}X_1 + \beta_2X_2
  + \beta_4\frac{X_3X_4}{\beta_3}
  + \cdots + \cos(\beta_pX_p).
\]
This is still parametric because it is controlled by a finite number of parameters, but it is not necessarily sensible, interpretable, or easy to fit.  Parametric modeling is powerful because it reduces an infinite function-estimation problem to a finite parameter-estimation problem.  Its weakness is that the chosen form may be wrong.

\subsection{Nonparametric Methods}

A method is nonparametric if it does not impose a fixed finite-dimensional form on $f$ in advance.  Nonparametric methods attempt to learn the shape of $f$ more directly from the data.  They can be deceptively simple or highly complex.

The advantage of a nonparametric method is flexibility.  If the true relationship is highly nonlinear, a nonparametric method may capture structure that a linear model misses.  The disadvantage is that flexibility requires data.  Without enough observations, a flexible method may not be able to distinguish signal from noise.  This is another expression of the bias--variance tradeoff: nonparametric methods often reduce bias but may increase variance, especially when the sample size is small.

This course will use nonparametric methods, but the early chapters focus on linear models because linear models make the statistical logic visible.  Understanding prediction error, bias, variance, residuals, and validation in simple models prepares us to use more complex models responsibly.

\section{Model Flexibility and Interpretability}

Different models occupy different positions on the flexibility--interpretability spectrum.  Least squares linear regression is relatively inflexible but highly interpretable.  Subset selection and lasso remain close to the linear-model family but modify which predictors are used or how coefficients are estimated.  Generalized additive models and trees are more flexible.  Bagging, boosting, support vector machines, and deep learning methods are often more flexible still, but their fitted relationships may be harder to interpret directly.

For prediction, increased flexibility can be valuable.  Modern prediction problems often involve nonlinear patterns, interactions, and high-dimensional structure that a basic linear model cannot capture.  For pure prediction tasks, linear models are often best treated as baselines rather than as the final word.

For inference, linear models remain extremely important.  Their assumptions, estimators, uncertainty measures, and diagnostic tools are well-developed.  When the goal is to explain relationships, defend assumptions, and communicate uncertainty to a decision maker, the interpretability of a linear model can be a major advantage.

The practical takeaway is not that one class of models is universally best.  The task determines the model.  If the goal is prediction, prioritize generalization error.  If the goal is inference, prioritize interpretability, assumptions, and uncertainty quantification.  When the goal includes both, state the tradeoff explicitly.

\section{Math Foundations Source Map}

This source map records the Math Foundations support used in this chapter and where those ideas appear in the completed Math Foundations textbook.

\begin{center}
\begin{tabular}{|p{0.24\textwidth}|p{0.36\textwidth}|p{0.28\textwidth}|}
\hline
\textbf{Chapter 2 use} & \textbf{Math Foundations connection} & \textbf{MF textbook location} \\
\hline
Response, prediction, and residual vectors & $\mathbb{R}^n$ vector notation, entries, and vector subtraction & Ch02 Secs. 2.1, 2.2, 2.6; Ch03 Sec. 3.1 \\
\hline
Empirical MSE as squared distance & Dot products, $\mathbf{a}^T\mathbf{a}=\sum_i a_i^2$, Euclidean norm, and distance between vectors & Ch05 Secs. 5.1, 5.3, 5.5, 5.7 \\
\hline
Residual length and RMSE & Euclidean norm properties and distance as $\|\mathbf{a}-\mathbf{b}\|$ & Ch02 Sec. 2.6; Ch05 Secs. 5.5, 5.7 \\
\hline
Linear models as finite-parameter estimators & Linear combinations, coefficients of linear combinations, and inner products & Ch03 Sec. 3.5; Ch05 Secs. 5.1, 5.3 \\
\hline
Flexibility and degrees-of-freedom intuition & Linear independence, bases, and the maximum size of linearly independent sets in $\mathbb{R}^n$ & Ch02 Sec. 2.1; Ch06 Secs. 6.3, 6.4, 6.5 \\
\hline
Matrix notation behind fitted predictors & Matrix size, transpose notation, and matrix-vector multiplication & Ch02 Sec. 2.6; Ch07 Secs. 7.1, 7.4, 7.6.1, 7.6.2 \\
\hline
\end{tabular}
\end{center}

\section{Chapter Summary}

This chapter developed the statistical meaning of prediction error.  Mean squared error measures average squared prediction error.  Root mean squared error places the score back in the units of the response.  The bias--variance decomposition shows that expected prediction error consists of squared bias, model variance, and irreducible noise.

Bias is systematic error caused by the model missing the true relationship on average.  Model variance is instability caused by sensitivity to the training data.  Noise is irreducible error in the response itself.  More flexible models tend to reduce bias but increase variance, which is why more flexibility does not automatically mean better prediction.

Training MSE measures performance on data used to fit the model.  Test MSE measures performance on unseen data.  Because prediction is about generalization, test error is the more important quantity.  Train/test splits provide a first method for estimating generalization performance.

Finally, the chapter distinguished parametric and nonparametric estimators.  Parametric methods assume a finite-parameter structure for $f$, while nonparametric methods avoid imposing a fixed finite-dimensional form.  Both can be useful, but both must be evaluated by whether they support the goal of the analysis.
}{
\section{Why Prediction Error Matters}

The previous chapter introduced the supervised learning model
\begin{equation}
  Y=f_0(\mathbf{X})+\epsilon,
  \label{eq:ch2-true-model}
\end{equation}
where $f_0$ is the true systematic relationship and $\epsilon$ is unexplained error.  This chapter asks a prediction-focused question:
\begin{quote}
  Given the information in $\mathbf{X}$, how accurately can we predict $Y$?
\end{quote}
We write $f_0$ for the true unknown function and $\hat f$ for a fitted model.  The fitted model is random because it depends on the training data.  A different training sample would usually produce a different $\hat f$.

\section{Mean Squared Error}

For a new observation with predictors $\mathbf{x}_0$ and response $Y_0$, the prediction is
\[
  \widehat{Y}_0=\hat f(\mathbf{x}_0).
\]
The mean squared error at $\mathbf{x}_0$ is
\begin{equation}
  \operatorname{MSE}(\mathbf{x}_0)
  =\mathbb{E}\left[\left(Y_0-\hat f(\mathbf{x}_0)\right)^2\right].
  \label{eq:ch2-mse-fixed-x}
\end{equation}
For observed data, the empirical MSE is
\begin{equation}
  \operatorname{MSE}
  =\frac{1}{n}\sum_{i=1}^n (y_i-\hat y_i)^2.
  \label{eq:ch2-empirical-mse}
\end{equation}
The root mean squared error is
\begin{equation}
  \operatorname{RMSE}=\sqrt{\operatorname{MSE}},
  \label{eq:ch2-rmse}
\end{equation}
which is useful because it is measured in the same units as the response.

\paragraph{Math Foundations Connection.}
Let $\mathbf{y}$ be the response vector, $\widehat{\mathbf{y}}$ the prediction vector, and $\mathbf{e}=\mathbf{y}-\widehat{\mathbf{y}}$ the residual vector.  Then
\begin{equation}
  \operatorname{MSE}=\frac{1}{n}\|\mathbf{y}-\widehat{\mathbf{y}}\|^2.
  \label{eq:ch2-mse-norm}
\end{equation}
Thus empirical MSE is a scaled squared distance.  See the Math Foundations textbook, Chapters 2, 3, 5, and 7.

\section{The Bias--Variance Decomposition}

For a fixed predictor value $\mathbf{x}_0$, suppose
\begin{equation}
  Y_0=f_0(\mathbf{x}_0)+\epsilon_0,
  \qquad
  \mathbb{E}[\epsilon_0]=0,
  \qquad
  \operatorname{Var}(\epsilon_0)=\sigma^2.
  \label{eq:ch2-new-response}
\end{equation}
Then the expected squared prediction error decomposes as
\begin{equation}
  \boxed{
  \mathbb{E}\left[\left(Y_0-\hat f(\mathbf{x}_0)\right)^2\right]
  =\operatorname{Bias}(\hat f(\mathbf{x}_0))^2
  +\operatorname{Var}(\hat f(\mathbf{x}_0))
  +\sigma^2.}
  \label{eq:ch2-bias-variance-final}
\end{equation}
The three terms have different meanings:
\begin{itemize}
  \item \textbf{Bias} is systematic error from using a model class that misses the true relationship on average.
  \item \textbf{Model variance} is how much the fitted model changes across different training samples.
  \item \textbf{Noise} is irreducible error that remains even if the systematic relationship is known perfectly.
\end{itemize}
More flexible models can reduce bias, but they usually increase variance.  Good prediction is a balance, not simply the most flexible fit.

\section{Training Error, Test Error, and Flexibility}

Training MSE measures performance on data used to fit the model:
\begin{equation}
  \operatorname{MSE}_{\text{train}}
  =\frac{1}{n_{\text{train}}}\sum_{(\mathbf{x}_i,y_i)\in\mathcal{D}_{\text{train}}}
  \left(y_i-\hat f(\mathbf{x}_i)\right)^2.
  \label{eq:ch2-training-mse}
\end{equation}
Test MSE measures performance on data not used to fit the model:
\begin{equation}
  \operatorname{MSE}_{\text{test}}
  =\frac{1}{n_{\text{test}}}\sum_{(\mathbf{x}_i,y_i)\in\mathcal{D}_{\text{test}}}
  \left(y_i-\hat f(\mathbf{x}_i)\right)^2.
  \label{eq:ch2-test-mse}
\end{equation}
For prediction, test error matters more because it measures generalization.

A model's \emph{flexibility} is its ability to represent a wide range of shapes.  As a useful rule of thumb,
\begin{equation}
  \text{more flexibility} \quad \Longrightarrow \quad
  \text{lower bias but higher variance}.
  \label{eq:ch2-flexibility-rule}
\end{equation}
This is why training MSE usually decreases as flexibility increases, while test MSE can decrease at first and then increase.

\section{Overfitting and Train/Test Splits}

\emph{Overfitting} occurs when a model follows noise or accidental patterns in the training data.  The model may look good on the training sample but perform poorly on new observations.  The goal is not to make training MSE as small as possible; the goal is to choose a model that generalizes.

A train/test split gives a first check on generalization:
\begin{enumerate}
  \item split the data into training and test sets;
  \item fit the model using only the training set;
  \item predict the test responses;
  \item compute test MSE or RMSE.
\end{enumerate}
The test set must remain unseen during fitting and model tuning.

The lecture simulation used data generated from
\begin{equation}
  Y=f_0(X)+\epsilon,
  \qquad
  f_0(X)=X^3,
  \qquad
  \epsilon\sim\mathcal{N}(0,9^2).
  \label{eq:ch2-simulation-model}
\end{equation}
A linear estimate underfit the cubic truth; a moderately flexible estimate did better; and an overly flexible estimate chased noise.  Training error alone rewards flexibility, but test error rewards generalization.

\section{Parametric and Nonparametric Estimators}

A \emph{parametric} method assumes a specific functional form with a finite set of parameters.  For example,
\begin{equation}
  f(\mathbf{X})=\beta_0+\beta_1X_1+\cdots+\beta_pX_p
  \label{eq:ch2-parametric-linear}
\end{equation}
reduces learning to estimating $\beta_0,\ldots,\beta_p$.  This is easier to fit and interpret, but it can be biased if the assumed form is too simple.  The expression above is a linear combination of input variables; see the Math Foundations textbook, Chapters 2, 3, 5, and 7.

A \emph{nonparametric} method avoids committing to a simple fixed form for $f$.  This can reduce bias, but it usually requires more data and can increase variance.  The practical question is which approach gives defensible prediction or inference for the task.

\section{Model Flexibility and Interpretability}

Flexible models can be powerful for prediction, but they may be harder to interpret.  Linear models are often less flexible, but their assumptions, coefficients, uncertainty measures, and diagnostics are easier to explain.  If the goal is prediction, prioritize generalization error.  If the goal is inference, prioritize interpretation, assumptions, and uncertainty quantification.  If the goal includes both, state the tradeoff explicitly.

\section{Math Foundations Source Map}

The Core Reader keeps the short Math Foundations connections for MSE as squared distance and parametric models as linear combinations.  The full Math Foundations textbook-location map for this chapter appears in the full reference textbook.

\section{Chapter Summary}

MSE measures average squared prediction error, and RMSE puts the score back in response units.  The bias--variance decomposition separates expected prediction error into squared bias, model variance, and irreducible noise.  Training error measures fit to data already seen; test error measures generalization.  More flexibility can reduce bias but increase variance, so the best predictive model is usually a balance between underfitting and overfitting.
}
